\begin{helper}
For every two-bite sample \(\omega\), the shadow graph of
\(\noderef{TwoBiteFinalGraph}\), obtained by forgetting colors of surviving
edges, has no triangle on three distinct vertices.
\end{helper}
\begin{proof}
Let \(G_R\), \(G_B\), and \(G_0\) be the final red graph, final blue graph, and
shadow graph.  Suppose that distinct vertices \(u,v,w\) form a triangle in
\(G_0\).  By the definition of the shadow graph in
\(\noderef{TwoBiteFinalGraph}\), each shadow edge is witnessed by a surviving
colored edge: for each of \(uv,uw,vw\), choose either a final red edge or a
final blue edge witnessing that shadow adjacency.

There are eight possible color choices for the three witnessed edges.  If all
three are red, the witnesses form a red triangle, contradicting
\(\noderef{TwoBiteFinalGraphNoRedTriangle}\).  If all three are blue, they form
a blue triangle, contradicting \(\noderef{TwoBiteFinalGraphNoBlueTriangle}\).

In every remaining case exactly two witnessed edges have one color and the
third has the other color.  The two same-colored edges share one of the
vertices \(u,v,w\).  If the repeated color is red, relabel that common endpoint
as the center; the two red edges from the center and the opposite blue edge
give the forbidden pattern of
\(\noderef{TwoBiteFinalGraphNoRedRedBlueTriangle}\).  If the repeated color is
blue, the same relabelling gives the forbidden pattern of
\(\noderef{TwoBiteFinalGraphNoBlueBlueRedTriangle}\).  Thus every possible
coloring of the three shadow edges contradicts one of the four colored
triangle exclusions, so \(G_0\) is triangle-free.
\end{proof}
